% Forms of parallel hardware, as a single nesting story:
%
%   Core  = A/F + two F/D (hyperthreads) + L1
%   CPU   = several Cores sharing an L2/L3
%   Node  = several CPUs sharing memory
%   Cluster = several Nodes connected by a network
%
% The five panels are the same structure at five levels of zoom, which is
% exactly why this is worth drawing with macros rather than by hand: \drawcore
% is written once and ends up on screen sixteen times in the last panel.
\documentclass[dvisvgm,border=3pt]{standalone}
\input{preamble.tex}
\begin{document}
\begin{tikzpicture}[
    x=1cm,y=1cm,
    box/.style ={chip,font=\sffamily\tiny},
    lab/.style ={text=inkdark,font=\sffamily\tiny,anchor=south west},
    up/.style  ={->,draw=uibknavy,line width=0.7pt,>={Stealth[round,length=2pt,width=2.4pt]}},
    pan/.style ={text=fg,font=\sffamily\small,align=center}]

% --------------------------------------------------------------------------
% one core: two hardware threads feeding one set of execution units
% --------------------------------------------------------------------------
\newcommand{\drawcore}[1]{%
  \begin{scope}[shift={#1}]
    \node[container,minimum width=15mm,minimum height=20mm,anchor=south west] at (0,0) {};
    \node[box,minimum width=13mm,minimum height=3.4mm] at (0.75,0.62) {L1 Cache};
    \node[box,minimum width=5.6mm,minimum height=3.2mm] at (0.42,1.15) {F/D};
    \node[box,minimum width=5.6mm,minimum height=3.2mm] at (1.08,1.15) {F/D};
    \node[box,minimum width=13mm,minimum height=3.4mm] at (0.75,1.68) {A/F};
    \draw[up] (0.42,0.79) -- (0.42,0.99);
    \draw[up] (1.08,0.79) -- (1.08,0.99);
    \draw[up] (0.42,1.31) -- (0.42,1.51);
    \draw[up] (1.08,1.31) -- (1.08,1.51);
    \node[lab] at (0.07,0.06) {Core};
  \end{scope}}

% --------------------------------------------------------------------------
% one CPU: two cores sharing an L2/L3
% --------------------------------------------------------------------------
\newcommand{\drawcpu}[1]{%
  \begin{scope}[shift={#1}]
    \node[container,minimum width=34mm,minimum height=36mm,anchor=south west] at (0,0) {};
    \drawcore{(0.12,1.35)}
    \drawcore{(1.78,1.35)}
    \node[box,minimum width=29mm,minimum height=4mm] at (1.70,0.95) {L2/L3 Cache};
    \draw[up] (0.87,1.15) -- (0.87,1.35);
    \draw[up] (2.53,1.15) -- (2.53,1.35);
    \node[lab] at (0.09,0.09) {CPU};
  \end{scope}}

% --------------------------------------------------------------------------
% one node: two CPUs sharing memory
% --------------------------------------------------------------------------
\newcommand{\drawnode}[1]{%
  \begin{scope}[shift={#1}]
    \node[nodebox,minimum width=72mm,minimum height=43mm,anchor=south west] at (0,0) {};
    \drawcpu{(0.15,0.45)}
    \drawcpu{(3.65,0.45)}
    \node[text=onfill,font=\sffamily\tiny,anchor=south west] at (0.12,0.12) {Node};
  \end{scope}}

% ==========================================================================
% panel 1 -- a plain single-core CPU
% ==========================================================================
\begin{scope}[xshift=0cm]
  \node[container,minimum width=17.5mm,minimum height=36mm,anchor=south west] at (0,0) {};
  \node[box,minimum width=15mm,minimum height=5mm] at (0.875,0.95) {L2/L3 Cache};
  \node[box,minimum width=15mm,minimum height=5mm] at (0.875,1.70) {L1 Cache};
  \node[box,minimum width=15mm,minimum height=5mm] at (0.875,2.45) {Fetch/Decode};
  \node[box,minimum width=15mm,minimum height=5mm] at (0.875,3.15) {ALU/FPU};
  \foreach \a/\b in {1.20/1.45, 1.95/2.20, 2.70/2.90}{\draw[up] (0.875,\a) -- (0.875,\b);}
  \node[lab] at (0.09,0.09) {CPU};
  \node[pan] at (0.875,-0.95) {single-core\\CPU};
\end{scope}

% ==========================================================================
% panel 2 -- the same core, now with two hardware threads
% ==========================================================================
\begin{scope}[xshift=2.45cm]
  \node[container,minimum width=19.5mm,minimum height=36mm,anchor=south west] at (0,0) {};
  \node[box,minimum width=17mm,minimum height=5mm] at (0.975,0.95) {L2/L3 Cache};
  \node[box,minimum width=17mm,minimum height=5mm] at (0.975,1.70) {L1 Cache};
  \node[box,minimum width=7.2mm,minimum height=5mm] at (0.60,2.45) {F/D};
  \node[box,minimum width=7.2mm,minimum height=5mm] at (1.35,2.45) {F/D};
  \node[box,minimum width=17mm,minimum height=5mm] at (0.975,3.15) {ALU/FPU};
  \draw[up] (0.975,1.20) -- (0.975,1.45);
  \draw[up] (0.60,1.95) -- (0.60,2.20);
  \draw[up] (1.35,1.95) -- (1.35,2.20);
  \draw[up] (0.60,2.70) -- (0.60,2.90);
  \draw[up] (1.35,2.70) -- (1.35,2.90);
  \node[lab] at (0.09,0.09) {CPU};
  \node[pan] at (0.975,-0.95) {single-core\\CPU + HT};
\end{scope}

% ==========================================================================
% panel 3 -- two such cores on one CPU
% ==========================================================================
\begin{scope}[xshift=5.1cm]
  \drawcpu{(0,0)}
  \node[pan] at (1.70,-0.95) {multi-core\\CPU + HT};
\end{scope}

% ==========================================================================
% panel 4 -- two CPUs in one shared-memory node
% ==========================================================================
\begin{scope}[xshift=9.2cm]
  \drawnode{(0,0)}
  \node[pan] at (3.60,-0.95) {shared memory node of\\multi-core CPUs + HT};
\end{scope}

% ==========================================================================
% panel 5 -- four such nodes, wired into a cluster
% ==========================================================================
\begin{scope}[xshift=17.3cm,yshift=1.36cm,scale=0.52,transform shape]
  \drawnode{(0,3.1)}
  \drawnode{(8.4,3.1)}
  \drawnode{(0,-2.6)}
  \drawnode{(8.4,-2.6)}
  % network between the four nodes
  \draw[wire,line width=1.6pt] (7.2,5.25) -- (8.4,5.25);
  \draw[wire,line width=1.6pt] (7.2,-0.45) -- (8.4,-0.45);
  \draw[wire,line width=1.6pt] (3.6,3.1) -- (3.6,1.7);
  \draw[wire,line width=1.6pt] (12.0,3.1) -- (12.0,1.7);
  \draw[wire,line width=1.6pt] (7.2,3.1) -- (8.4,1.7);
  \draw[wire,line width=1.6pt] (8.4,3.1) -- (7.2,1.7);
\end{scope}
\node[pan] at (21.4,-0.95) {cluster of shared memory\\nodes of multi-core CPUs + HT};

\end{tikzpicture}
\end{document}
